\documentclass[12pt]{article}
\usepackage[margin=1in]{geometry}
\usepackage{amsmath}
\usepackage{fancyhdr}
\fancypagestyle{plain}{%
  \fancyhf{}
  \fancyhead[L]{POLISCI 450A}
  \fancyhead[R]{Fall 2026}
  %\fancyfoot[C]{\thepage}
}

\usepackage{titlesec}
\titleformat{\section}{\large\bfseries}{\thesection}{1em}{}
\usepackage{titling}
\setlength{\droptitle}{-3em}   % move title up
\usepackage{setspace}

\usepackage{amsmath, amssymb, amsthm, graphicx, xcolor, bm}
\newcommand{\bs}[1]{\bm{#1}}
\newcommand{\var}{\operatorname{Var}}
\newcommand{\cov}{\operatorname{Cov}}
\newcommand{\E}{\operatorname{\mathbb{E}}}
\newenvironment{solution}{\par\noindent\textbf{Solution.}\ }{\hfill$\square$\par}

\DeclareMathOperator*\plim{plim}
\DeclareMathOperator{\cor}{Cor}
\newcommand{\indep}{\perp\!\!\!\perp}
\newcommand{\cN}{\mathcal{N}}

\providecommand{\mb}{\mathbf}
\newcommand{\X}{\ensuremath{\mathbf{X}}}
\newcommand{\y}{\ensuremath{\mathbf{y}}}
\newcommand{\M}{\ensuremath{\mathbf{M}}}
\renewcommand{\P}{\ensuremath{\mathbf{P}}}
\renewcommand{\u}{\ensuremath{\mathbf{u}}}  % override accent for vector u
\newcommand{\ehat}{\hat{\boldsymbol{\epsilon}}}
\newcommand{\e}{\ensuremath{\bs{\epsilon}}}
\newcommand{\I}{\ensuremath{\mathbf{I}}}
\newcommand{\XX}{\ensuremath{\X' \X}}
\newcommand{\XXinv}{\ensuremath{(\X' \X)^{-1}}}
\newcommand{\hatmat}{\ensuremath{\X \XXinv \X'}}
\renewcommand{\a}{\mathbf{a}}
\renewcommand{\b}{\ensuremath{\bs{\beta}}}
\newcommand{\bhat}{\hat{\boldsymbol{\beta}}}

\newcommand{\hatsig}{\ensuremath{\hat{\sigma}^2}}
\newcommand{\real}{\mathbb{R}}
\newcommand{\R}{\mathbf{R}}
\newcommand{\A}{\mathbf{A}}
\renewcommand{\r}{\mathbf{r}}

\title{\Large Lecture 6-2: OLS Inference}
\author{}
\date{}
\begin{document}
\maketitle
\vspace{-6em}

\noindent \textbf{Goal:} Develop statistical inference methods for multiple regression coefficients.\medskip

\begin{itemize}\setlength{\itemsep}{1em}
\item Testing hypotheses about individual coefficients
\begin{itemize}
\item $t$-values in multiple regression
\item Test procedure
\item Confidence intervals
\end{itemize}
\item Testing linear hypotheses: simple case
\begin{itemize}
\item Linear combinations of coefficients
\item Variance of linear combinations
\end{itemize}
\item Testing joint significance
\begin{itemize}
\item $F$ tests for joint hypotheses
\item $\chi^2$ and $F$ distributions
\item Restricted vs unrestricted models
\item Ex1. Omnibus tests
\item Ex2. Equality of coefficients tests
\item Relationship between $t$ and $F$ tests
\end{itemize}
\item Testing linear hypotheses: general case
\begin{itemize}
\item Limitations of simple $F$ tests
\item General restriction matrix form
\item Wald statistic
\item Large and small sample distributions
\end{itemize}
\end{itemize}

% End outline
\clearpage


\section*{Worksheet}


\begin{enumerate}\itemsep1em


% \item (\emph{t–statistic and its small–sample distribution in multiple regression})  
% Consider the homoskedastic Gauss–Markov model with normal errors
% \[
% \y=\X\b+\e,
% \qquad \e\sim\cN(0,\sigma^2\I_n),
% \quad \text{rank}(\X)=k+1<n.
% \]
% Let $\hat\b=(\X'\X)^{-1}\X'\y$ and $\hatsig=\frac{1}{n-k-1}\,\ehat'\ehat$ with $\ehat=\y-\X\hat\b$. Fix $j\in\{0,\ldots,k\}$ and a scalar $c$. \\ \bigskip

% Explain why the test statistic $T\;=\;\frac{\hat\beta_j-c}{\widehat{SE}(\hat\beta_j)}$, where  $\widehat{SE}(\hat\beta_j)=\sqrt{\hatsig\;\big[(\X'\X)^{-1}\big]_{(j,j)}}$, has the $t_{\,n-k-1}$ distribution under $H_0:\beta_j=c$.\\ 

% \textbf{Hint}: Need to show (a) $\hat\beta_j$ is normally distributed; (b) $\frac{(n-k-1)\,\hatsig}{\sigma^2}$ is $\chi^2_{,n-k-1}$ distributed; (c) the two parts are independent; and (d) normal-$\chi^2$ ratio characterization. Explain why. 

% \clearpage

\item  (\emph{A linear combination $(\beta_{PS}-2\beta_{S})$ and its standard error}) Return to the Chile model with dummies $S$ (secondary) and $PS$ (post-secondary), with primary schooling as the baseline. The estimated model is
\[
\widehat{Y} = \hat\beta_0 + \hat\beta_{fem}\,FEM + \hat\beta_{age}\,AGE + \hat\beta_{S}\,S + \hat\beta_{PS}\,PS.
\]
We wish to test whether the 1-2-3 coding of education fits the data:
\[
H_0:\ \beta_{PS}-2\beta_{S}=0.
\]
\begin{enumerate}
\item Derive $\var(\hat\beta_{PS}-2\hat\beta_{S})$
\item Let $\mb{A}=(\X'\X)^{-1}=\big[a_{rs}\big]$ and $\widehat{\var}(\hat\b)=\hatsig\,\mb{A}$. Explain that
\[
\widehat{\var}\!\big(\hat\beta_{PS}-2\hat\beta_{S}\big)
= \hatsig\Big(a_{PS,PS}+4\,a_{S,S}-4\,a_{S,PS}\Big)
\]
is an unbiased estimator of $\var(\hat\beta_{PS}-2\hat\beta_{S}\mid\X)$ under Assumptions I--V.
\end{enumerate}

\clearpage

\item Consider the multiple regression model
\[
Y = \beta_0 + \beta_1 X_1 + \beta_2 X_2 + \cdots + \beta_k X_k + u,
\]
and suppose we wish to test the joint null hypothesis
\[
H_0:\ \beta_1=\beta_2=\cdots=\beta_k=0
\quad\text{versus}\quad
H_1:\ \text{at least one }\beta_j\neq 0.
\]

\begin{enumerate}
\item Explain in words what this hypothesis is asking about the explanatory power of the regressors.
\item Describe step-by-step how to conduct the $F$ test for this hypothesis, including:
  \begin{enumerate}
    \item the unrestricted and restricted models,  
    \item how to compute the $F$ statistic from the two sum of squared residuals, and  
    \item how to determine whether to reject $H_0$ at significance level $\alpha$.
  \end{enumerate}
\item Interpret the meaning of rejecting versus not rejecting $H_0$ in terms of the overall explanatory power of $X_1,\ldots,X_k$.
\end{enumerate}

\clearpage
\item Consider the multiple regression model
\[
Y = \beta_0 + \beta_1 X_1 + \beta_2 X_2 + \beta_3 X_3 + \cdots + \beta_k X_k + u,
\]
and suppose we want to test whether the first three regressors have the same effect:
\[
H_0:\ \beta_1=\beta_2=\beta_3 \quad \text{versus} \quad H_1:\ \text{at least one differs.}
\]

\begin{enumerate}
% \item Express $H_0$ as a system of linear restrictions in matrix form, and determine how many independent restrictions are being tested.
\item Describe the step-by-step $F$ test procedure, including:
  \begin{enumerate}
    \item how to construct the restricted model implied by $H_0$,
    \item how to compute the $F$ statistic using the two sums of squared residuals, and
    \item the reference distribution used for inference.
  \end{enumerate}
\item Provide an intuitive interpretation of what rejecting or failing to reject $H_0$ implies about the relationship between $X_1$, $X_2$, and $X_3$.
\end{enumerate}
\clearpage

\item \textit{(General procedure for testing linear hypotheses)} Consider the linear regression model
\[
Y = \beta_0 + \beta_1 X_1 + \beta_2 X_2 + \beta_3 X_3 + u,
\]
with coefficient vector $\b = [\beta_0\ \beta_1\ \beta_2\ \beta_3]'$. We wish to test linear hypotheses of the form $R\b=r$. Construct the corresponding $R$ and $r$ in the following two scenarios.
\begin{enumerate}
\item Test the hypothesis that all three slope coefficients equal 3:
\[
H_0:\ \beta_1=\beta_2=\beta_3=3.
\]
Write this null hypothesis in matrix form $R\b=r$ by explicitly constructing $R$ and $r$.

\item Test the composite hypothesis
\[
H_0:\ \beta_1 = 2\beta_2 = 0.5\beta_3 + 1.
\]
\end{enumerate}

\end{enumerate}


\end{document}
